\documentclass[11pt,a4paper]{article}
\usepackage[T1]{fontenc}
\usepackage[utf8]{inputenc}
\usepackage{lmodern}
\usepackage[margin=27mm]{geometry}
\usepackage{amsmath,amssymb,amsthm,mathtools}
\usepackage{enumitem,booktabs,array}
\usepackage{xurl}
\usepackage[colorlinks=true,linkcolor=blue,urlcolor=blue]{hyperref}
\usepackage{microtype}
\setlength{\emergencystretch}{2em}
\setcounter{tocdepth}{2}
\newtheorem{theorem}{Theorem}[section]
\newtheorem{lemma}[theorem]{Lemma}
\newtheorem{proposition}[theorem]{Proposition}
\newtheorem{corollary}[theorem]{Corollary}
\theoremstyle{definition}
\newtheorem{definition}[theorem]{Definition}
\newtheorem{obligation}[theorem]{Construction obligation}
\newtheorem{example}[theorem]{Example}
\theoremstyle{remark}
\newtheorem*{remark}{Remark}
\newcommand{\N}{\mathbb N}
\newcommand{\Z}{\mathbb Z}
\newcommand{\R}{\mathbb R}
\newcommand{\Co}{C_{\mathrm{out}}}
\newcommand{\Cs}{C_{\mathrm{sc}}}
\newcommand{\tot}[1]{\operatorname{tot}(#1)}
\newcommand{\Reach}{\operatorname{Reach}}
\newcommand{\Hit}{\operatorname{Hit}}
\title{Symmetric Three-Counter Reachability:\\
Precise Definitions, Proof Ideas, and the Remaining Construction}
\author{}
\date{1 November 2026}
\begin{document}
\maketitle

\begin{abstract}
The proposed decision procedure separates two tasks. The first is
geometric: replace machine runs by compositions of explicitly specified
relations. The second is algorithmic: decide those compositions using
finite bases of upward sets. This text defines the objects used in both
tasks, explains the purpose of each definition, and proves the
finite-graph and saturation arguments in detail.
It also corrects the status of the earlier version. The effective
construction of one fixed family of relations, preserved by all the
proposed replacements of runs, has not been fully established.
The precise remaining obligation is stated below; the two assumed
semilinearity results alone are not presented as a completed proof of
three-counter decidability.
\end{abstract}

\paragraph{How to read this version.}
Section~\ref{sec:model} starts with the actual machine and gives an exact
boundary decomposition. This shows where the difficulty lies before any
auxiliary terminology is introduced.
Section~\ref{sec:formal-saturation} proves the algorithms for composing
relations with pumps.
Sections~\ref{sec:pair-rigorous} and \ref{sec:face-query} explain how a
bound on two physical counters produces the finite descriptions needed
by those algorithms.
Sections~\ref{sec:records}--\ref{sec:construction} state precisely what
must still be constructed from the original machine.
An informal name is never intended to supply a missing hypothesis:
for example, calling a relation a positive passage does not prove
its pumping or monotonicity properties.

\tableofcontents
\bigskip

\section{The machine, the assumptions, and the actual difficulty}
\label{sec:model}

\subsection{What symmetry means}

A counter stores a nonnegative integer. An update changes several
counters at once, but is legal only if no counter becomes negative.
A zero test changes the control state without changing the counters.
Symmetry means that an available update can be applied with any
permutation of its coordinate changes. It does not mean that a run can
be reversed.

\begin{definition}[Counter machines and coordinate symmetry]
\label{def:machine}
For \(d\ge1\), a \(d\)-counter machine is a triple
\[
 M=(Q,\Delta,Z),
 \qquad
 \Delta\subseteq Q\times\Z^d\times Q,\quad
 Z\subseteq Q\times\{1,\ldots,d\}\times Q,
\]
where all three sets are finite. Its configurations form
\(\operatorname{Conf}(M)=Q\times\N^d\), with \(0\in\N\).
Its step relation is the union of the following relations:
\[
\begin{split}
 (p,x)\longrightarrow_M(q,x+\delta)
 &\quad\text{if }(p,\delta,q)\in\Delta
                 \text{ and }x+\delta\in\N^d,\\
 (p,x)\longrightarrow_M(q,x)
 &\quad\text{if }(p,i,q)\in Z\text{ and }x_i=0.
\end{split}
\]
Write \(\longrightarrow_M^*\) for reflexive-transitive closure and
\(\Reach(M)\) for this binary relation. A VASS is a machine with
\(Z=\varnothing\).

For a permutation \(\pi\) of \(\{1,\ldots,d\}\), set
\((\pi x)_i=x_{\pi^{-1}(i)}\). The machine is \emph{symmetric} if
\[
 (p,\delta,q)\in\Delta
 \quad\Longrightarrow\quad
 (p,\pi\delta,q)\in\Delta
\]
for every such permutation. Zero tests retain their specified indices.
\end{definition}

For example, an update \((-1,2,0)\) entails the updates obtained by
permuting these three entries, with the same source and target controls.
It does not entail the reverse update \((1,-2,0)\).
If a reversed machine is used in a proof, it is a different machine:
reverse the control edges and negate update vectors; reverse zero-test
control edges and keep their tested coordinate. Its runs are exactly
reversals of original runs, and it remains symmetric.

\subsection{What the two assumptions supply}

The assumed algorithms give formulas describing all pairs of endpoints
of certain machines. Such formulas can be composed a fixed finite
number of times by introducing existentially quantified intermediate
vectors. They do not automatically describe an unbounded number of
compositions.

\begin{definition}[Effective semilinearity]
\label{def:presburger}
A linear subset of \(\N^k\) is a set
\[
 b+\N p_1+\cdots+\N p_r,
 \qquad b,p_1,\ldots,p_r\in\N^k.
\]
A semilinear set is a finite union of linear sets.
Equivalently, it is definable in first-order Presburger arithmetic,
the structure \((\N,0,1,+,\le)\).
Effective semilinearity means that an algorithm outputs a finite
semilinear description, or equivalently a Presburger formula.
For finitely many control labels, a relation is effectively semilinear
when such a formula is computed for every pair of labels.
\end{definition}

The hypotheses under investigation are:
\begin{description}[style=nextline,leftmargin=2.8em]
\item[\textbf{S2}] Given a symmetric two-counter machine \(M\), compute
formulas for
\(\{(x,y):(p,x)\longrightarrow_M^*(q,y)\}\), for all \(p,q\).
\item[\textbf{S3}] Given a symmetric three-dimensional VASS \(V\),
compute the analogous formulas for all its control pairs.
\end{description}
The intended conclusion is an algorithm deciding, for an arbitrary
symmetric three-counter machine \(M\) and configurations \(c,d\),
whether \(c\longrightarrow_M^*d\).
That conclusion remains the goal; the completed implications proved
in this document are stated with their exact hypotheses.

For relations \(R\subseteq A\times B\) and \(S\subseteq B\times C\),
we use left-to-right composition:
\[
 R;S=\{(a,c):\exists b\in B,\ (a,b)\in R,\ (b,c)\in S\}.
\]
For \(I\subseteq A\) and \(J\subseteq B\),
\[
 I;R=\{b:\exists a\in I,\ aRb\},\qquad
 R^{-1}(J)=\{a:\exists b\in J,\ aRb\}.
\]
For a relation on one set, \(R^0\) is the identity relation,
\(R^{n+1}=R^n;R\), and \(R^*=\bigcup_{n\ge0}R^n\).
Finite union, intersection, complement, composition, and projection
of Presburger relations are effective. An expression \(R^*\) is
\emph{not} on this list.

\subsection{An exact reduction to relations on boundary faces}

Every zero test takes place where at least one counter is zero.
At such a configuration, only two counter values need to be stored.
We can therefore describe a run by its update-only pieces and its
zero-test steps, recording their boundary endpoints with two numbers.
This reduction is exact and uses no normalization assumption.

\begin{definition}[Zero-face records]
\label{def:zero-face}
For \(i\in\{1,2,3\}\), let
\(\iota_i:\N^2\to\N^3\) insert a zero at coordinate \(i\), putting the
two given values at the other coordinates in increasing index order.
Set
\[
 Y=\biguplus_{(q,i)\in Q\times\{1,2,3\}}
       \{(q,i)\}\times\N^2.
\]
The record \(u=(q,i,z)\) denotes the physical configuration
\(\rho(u)=(q,\iota_i(z))\). A configuration with two zeros can have
more than one record.

Let \(V=(Q,\Delta,\varnothing)\) be the update-only part of \(M\).
Define relations on \(Y\) by
\[
\begin{split}
 uA_{\mathrm{upd}}v
 &\iff \rho(u)\longrightarrow_V^*\rho(v),\\
 uA_{\mathrm{test}}v
 &\iff \rho(u)\longrightarrow_M\rho(v)
       \text{ by one zero-test transition},\\
 A&=A_{\mathrm{upd}}\cup A_{\mathrm{test}}.
\end{split}
\]
For fixed physical endpoints \(c,d\), define
\[
 I_c=\{u\in Y:c\longrightarrow_V^*\rho(u)\},
 \qquad
 J_d=\{u\in Y:\rho(u)\longrightarrow_V^*d\}.
\]
\end{definition}

\begin{proposition}[The exact boundary problem]
\label{prop:boundary}
Under \textbf{S3}, the relations \(A_{\mathrm{upd}},A_{\mathrm{test}}\)
and the sets \(I_c,J_d\) are effectively semilinear. Moreover
\begin{equation}
 c\longrightarrow_M^*d
 \quad\Longleftrightarrow\quad
 c\longrightarrow_V^*d
 \ \lor\
 (I_c;A^*)\cap J_d\ne\varnothing.
 \label{eq:raw-boundary}
\end{equation}
\end{proposition}
\begin{proof}
Apply \textbf{S3} to \(V\) and substitute zero in the recorded
coordinates to obtain \(A_{\mathrm{upd}},I_c,J_d\).
For \(A_{\mathrm{test}}\), take the finite disjunction over
\((p,k,q)\in Z\) requiring source control \(p\), target control \(q\),
equality of the physical vectors, and their \(k\)-th coordinate zero.
These are Presburger constraints.

If a run has no zero test, it is a \(V\)-run. Otherwise split it before
and after each zero test. Every endpoint at a split has a zero
coordinate and hence a record in \(Y\). Each update-only portion gives
an \(A_{\mathrm{upd}}\)-step, and each test gives an
\(A_{\mathrm{test}}\)-step. The first and last update-only portions
give membership in \(I_c\) and \(J_d\).
Conversely, each factor in such a composition has a realizing machine
run, and equal records at joins denote equal physical configurations.
Concatenation gives the required \(M\)-run.
\end{proof}

Thus the first real obstacle is already visible:
\(\,A\subseteq Y^2\,\) is a computable relation on a finite union of
copies of \(\N^2\), but its closure still has to be decided.
Merely calling \(A\) a collection of segments does not solve this.

For comparison, the one-dimensional relation
\[
 R=\{(n,2n):n\in\N\}
\]
is Presburger-definable, but \(\{1\};R^*=\{2^k:k\in\N\}\) is not
semilinear. Indeed every unary semilinear set is ultimately periodic,
whereas consecutive powers of two have unbounded gaps.
Also, no fixed \(D>0\) can make \(nR(2n)\) imply \(nR(2n+D)\), even
when the gain \(n\) is arbitrarily large.
This example is not claimed to be a symmetric-machine boundary
relation. It identifies what cannot follow from semilinearity alone.

\subsection{The main idea, before the technical definitions}

The proposed improvement of Proposition~\ref{prop:boundary} has three
parts.
\begin{enumerate}[leftmargin=*]
\item Some actual run families permit independent increases or
decreases of coordinates. Treat those operations separately.
Their images or preimages are upward sets, and upward sets have
finite bases.
\item Between those operations, a suitable numerical function should
increase at every remaining large excursion. If two endpoint
coordinates are bounded, that bounds the number of large excursions,
even though the third coordinate is unbounded.
\item A configuration with two bounded coordinates can be encoded by
one unbounded integer and finite control. Thus repeated visits to
such configurations lead to a scalar problem, where bounded gain
means finitely many possible differences. This is where \textbf{S2}
is intended to enter.
\end{enumerate}
The next sections make these assertions into precise mathematical
statements. They also separate their proofs from the construction
needed to apply them to every original machine.



% This section uses theorem, lemma, and definition environments.
% It assumes \N denotes the nonnegative integers.
\section{Why the two saturation procedures terminate}
\label{sec:formal-saturation}

The following argument concerns relations on tuples of natural numbers.
Its purpose is to isolate exactly what is needed from the preceding
geometric construction. In particular, it does not assume that an
arbitrary transitive closure of a semilinear relation is semilinear.

\subsection{The sets on which the algorithm operates}

A finite control label tells us which quantities are currently stored
as counters. Different labels may use different numbers of counters.
For example, a physical boundary face uses two counters, whereas an
axis with its two bounded coordinates recorded in the label uses one.
Keeping the label is essential: the same numerical tuple at two
different labels represents two different configurations.

\begin{definition}[Finite union of counter spaces]
\label{def:counter-spaces}
Fix a finite set \(Q\) and an integer \(d(q)\geq0\) for each \(q\in Q\).
Put
\[
 X=\biguplus_{q\in Q} X_q,
 \qquad X_q=\{q\}\times\N^{d(q)}.
\]
The disjoint union symbol means that the label \(q\) is part of the
element. We use the convention \(\N^0=\{()\}\), so a label of dimension
zero represents a single element.
\end{definition}

A pump may add a fixed period \(D\), rather than one. It then preserves
each counter's residue modulo \(D\). The appropriate order therefore
compares configurations only when both their control labels and their
residue vectors agree.

\begin{definition}[Grid order and upward closure]
\label{def:grid-order}
Fix an integer \(D>0\). For \(x=(q,\boldsymbol x)\) and
\(y=(q',\boldsymbol y)\) in \(X\), write \(x\preceq_D y\) if
\[
 q=q'
 \quad\text{and}\quad
 \boldsymbol y-\boldsymbol x\in D\N^{d(q)}.
\]
For \(A\subseteq X\), let
\[
 \mathord{\uparrow_D}A
   =\{y\in X:\exists x\in A,\ x\preceq_D y\}.
\]
A set \(A\) is \(D\)-upward if
\(\mathord{\uparrow_D}A=A\). A finite set \(F\) is a basis for \(A\)
if \(A=\mathord{\uparrow_D}F\); the basis need not be minimal.
\end{definition}

For each fixed label and residue vector, division by \(D\), after
subtracting that residue, identifies the order with the usual
coordinatewise order on \(\N^{d(q)}\). Dickson's lemma consequently
applies to \((X,\preceq_D)\). It gives two facts we will use:
every upward set has a finite basis, and every increasing sequence of
upward sets eventually becomes constant. For the second fact, if
\(A_0\subsetneq A_1\subsetneq\cdots\), choose
\(x_i\in A_{i+1}\setminus A_i\). Dickson's lemma gives \(i<j\) with
\(x_i\preceq_D x_j\), but then \(x_j\in A_j\), a contradiction.

The algorithm must also represent a singleton target, a bounded
interval, and the complement of its current approximation.
Products of intervals in fixed residues provide one effective family
containing all these sets.

\begin{definition}[Grid boxes]
\label{def:grid-boxes}
A grid interval is a set
\[
 I_D(r,\ell,u)
   =\{n\in\N:n\equiv r\pmod D,\ \ell\leq n\leq u\},
\]
where \(0\leq r<D\), \(\ell\in\N\), and
\(u\in\N\cup\{\infty\}\). The inequality \(n\leq\infty\) imposes no
restriction. A grid box in \(X_q\) is
\[
 \{q\}\times\prod_{i=1}^{d(q)} I_D(r_i,\ell_i,u_i).
\]
Let \(\mathcal B_D\) be the family of finite unions of grid boxes,
represented by their finite lists of bounds, residues, and labels.
\end{definition}

The family \(\mathcal B_D\) is effectively closed under union,
intersection, and complement relative to \(X\). Indeed, after splitting
by the finitely many labels and residues, these are Boolean operations
on finite unions of ordinary integer boxes. Their boundaries give a
finite common partition on which membership in each operand is constant.
Membership, emptiness, and equality are decidable from this
representation. Every singleton and every upward set given by a finite
basis belongs to \(\mathcal B_D\).

\subsection{From an intersection test to a finite basis}

Dickson's lemma asserts that a basis exists; it does not tell us when
an enumeration of reachable points has found the whole basis. The
extra information that solves this problem is the ability to ask
whether anything remains outside the approximation already computed.

\begin{lemma}[Effective basis computation]
\label{lem:effective-basis}
Let \(Y\subseteq X\) be \(D\)-upward. Suppose an algorithm decides,
for every \(E\in\mathcal B_D\), whether \(Y\cap E\neq\varnothing\).
Then a finite basis for \(Y\) can be computed.
\end{lemma}
\begin{proof}
Maintain a finite set \(F\) with
\(A=\mathord{\uparrow_D}F\subseteq Y\), initially \(F=\varnothing\).
Ask whether \(Y\cap(X\setminus A)\) is empty. If it is empty, output
\(F\). Otherwise enumerate the elements of the effectively enumerable
set \(X\), testing both membership in \(A\) and the singleton query
\(Y\cap\{x\}\neq\varnothing\), until a point \(x\in Y\setminus A\)
is found. Such a point exists by the preceding query. Add \(x\) to
\(F\), and repeat.

Since \(Y\) is upward, the new principal upward set generated by \(x\)
is contained in \(Y\). The approximation strictly increases at every
unsuccessful stopping test. The increasing-chain property proved
above implies termination.
\end{proof}

\subsection{The precise pumping relations}

An upward event lets each coordinate increase independently. This
means that \((x_1,x_2)\) can become
\((x_1+kD,x_2+\ell D)\) with unrelated \(k,\ell\in\N\).
It is stronger than the ability to add only \((D,D)\).
A downward event has the reversed property. Some events require a
fixed lower enabling buffer; an upward eligibility set records that
requirement.

\begin{definition}[Independent upward and downward events]
\label{def:independent-events}
An event description is a triple \((p,q,G)\), where
\(p,q\in Q\), \(d(p)=d(q)=d\), and \(G\subseteq\N^d\) is upward
for the order
\(\boldsymbol x\preceq_D\boldsymbol y
\iff\boldsymbol y-\boldsymbol x\in D\N^d\).
The set \(G\) is supplied by a finite basis.
Its upward and downward relations are, respectively,
\[
\begin{split}
 U_{p,q,G}
 &=\{((p,\boldsymbol x),(q,\boldsymbol y)):
       \boldsymbol x,\boldsymbol y\in G,\
       \boldsymbol y-\boldsymbol x\in D\N^d\},\\
 V_{p,q,G}
 &=\{((p,\boldsymbol x),(q,\boldsymbol y)):
       \boldsymbol x,\boldsymbol y\in G,\
       \boldsymbol x-\boldsymbol y\in D\N^d\}.
\end{split}
\]
Let \(U\) be a finite union of specified upward event relations and
\(V\) a finite union of specified downward event relations.
The input and output labels may differ, and zero change is permitted.
\end{definition}

Every image \(A;U\), for any \(A\subseteq X\), is upward. Every
preimage \(V^{-1}(A)\) is upward. Both assertions concern the labelled
space \(X\). For example, if \(xVy\) and \(x\preceq_D x'\), then
\(x'Vy\): the larger input remains eligible and can subtract to the
same output.

For represented box sets \(E\), the sets \(E;V\) and \(U^{-1}(E)\)
are effectively in \(\mathcal B_D\). To check this, consider one
event and rename its input and output labels as necessary. Write
\[
 \mathord{\downarrow_D}A
  =\{\boldsymbol x:\exists\boldsymbol y\in A,\
                         \boldsymbol x\preceq_D\boldsymbol y\}.
\]
Then the numerical part of \(E;V_{p,q,G}\) is
\(G\cap\mathord{\downarrow_D}(E_p\cap G)\), where
\(E_p=\{\boldsymbol x:(p,\boldsymbol x)\in E\}\).
The numerical part of \(U_{p,q,G}^{-1}(E)\) is
\(G\cap\mathord{\downarrow_D}(E_q\cap G)\).
Discard empty boxes first. Taking downward closure of a nonempty box removes its lower bounds
within its residue class, so both expressions are effective box sets.

\subsection{Deciding alternating passages and pumping events}

Between two events there is a relation \(K\). The algorithm will not
need a formula for all pairs in \(K\), nor a formula for \(K^*\).
It only needs a test for whether a \(K\)-passage can connect two
represented box sets. This is a precise requirement on an algorithm,
not a claim that \(K\) is an ordinary machine transition.

\begin{definition}[Box-intersection algorithm for a fixed relation]
\label{def:box-intersection}
A box-intersection algorithm for \(K\subseteq X\times X\) decides,
given \(S,T\in\mathcal B_D\), the proposition
\[
 \operatorname{Hit}_K(S,T)
 \quad\Longleftrightarrow\quad
 \exists x\in S\ \exists y\in T:\ (x,y)\in K.
\]
The same relation \(K\) is used for all inputs \(S,T\).
\end{definition}

For relations \(R,S\), composition is from left to right:
\(x(R;S)z\) means \(\exists y\,(xRy\land ySz)\).
Write \(A;R=\{y:\exists x\in A,\ xRy\}\) and
\(R^{-1}(A)=\{x:\exists y\in A,\ xRy\}\).
Reflexive-transitive closure includes the empty sequence of steps.

\begin{theorem}[Completion by independent events]
\label{thm:formal-completion}
Assume that \(U,V\) are given as in
Definition~\ref{def:independent-events} and that a box-intersection
algorithm for a fixed \(K\subseteq X\times X\) is supplied.
Then one can decide, for any \(S,T\in\mathcal B_D\), whether
\[
 (S;L)\cap T\neq\varnothing,
 \qquad
 L=K;\bigl((U\cup V);K\bigr)^*.
\]
No transitivity, reflexivity, or semilinearity assumption on \(K\)
is required.
\end{theorem}

\begin{proof}
We first close backwards under downward events, then close forwards
under upward events. Both operations become upward immediately after
the respective event, which is why increasing-chain termination
applies.

\emph{Backward computation.}
Put \(C=(V;K)^*\). For a represented target \(T\in\mathcal B_D\), set
\[
 T_0=T,\qquad
 T_{n+1}=T\cup V^{-1}\bigl(K^{-1}(T_n)\bigr).
 \tag{B}
\]
We explain how to compute each term without computing \(K^{-1}(T_n)\).
Let \(Y_n=V^{-1}(K^{-1}(T_n))\). It is upward, and, for any
\(E\in\mathcal B_D\),
\[
 Y_n\cap E\neq\varnothing
 \quad\Longleftrightarrow\quad
 \operatorname{Hit}_K(E;V,T_n).
\]
The right side is decidable because \(E;V\) is an effective box set.
Lemma~\ref{lem:effective-basis} computes \(Y_n\), so it computes
\(T_{n+1}\).

Inductively \(T_n\subseteq T_{n+1}\), hence \(Y_n\subseteq Y_{n+1}\).
The sets \(Y_n\) are upward, so they eventually stabilize. Equality of
successive \(T_n\)'s is decidable and is a stopping test. By induction,
\(T_n\) is exactly the set of points that reach \(T\) using at most
\(n\) occurrences of \(V;K\). The resulting fixed point is therefore
\(C^{-1}(T)\). In particular, for each given box set \(T\), this
preimage is effectively a box set; we have not computed the full
binary relation \(C\).

\emph{Forward computation.}
For \(A\in\mathcal B_D\), put
\[
 F(A)=A;K;C;U.
\]
The final upward event makes \(F(A)\) upward. For any
\(E\in\mathcal B_D\),
\[
 F(A)\cap E\neq\varnothing
 \quad\Longleftrightarrow\quad
 \operatorname{Hit}_K
 \bigl(A,C^{-1}(U^{-1}(E))\bigr).
\]
First compute the box set \(U^{-1}(E)\), then its \(C\)-preimage by
(B), then use the given algorithm for \(K\). Thus
Lemma~\ref{lem:effective-basis} computes a basis for \(F(A)\).

Starting from the given source \(S\), iterate
\[
 S_0=S,\qquad S_{n+1}=S\cup F(S_n).
 \tag{F}
\]
The sequence \(F(S_n)\) is increasing and upward, hence stabilizes.
Equality is decidable. Its resulting source set is
\(S;A_+\), where \(A_+=(K;C;U)^*\).

\emph{Final test.}
Split an alternating word in \(L\) immediately after each \(U\).
Each completed block is \(K;C;U\); the final block is \(K;C\).
Conversely every such concatenation is a word in \(L\). Therefore
\[
 L=A_+;K;C.
\]
This is an identity of relation expressions; it does not combine
adjacent \(K\)'s or require \(K\) to be transitive.
The desired answer is precisely
\[
 \operatorname{Hit}_K\bigl(S;A_+,C^{-1}(T)\bigr).
\]
Both arguments have already been computed as box sets.
\end{proof}

\subsection{The one-counter version used between axis visits}

In the scalar application, some transitions already have Presburger
formulas. A positive transition may increase the output by \(D\)
while keeping its input fixed. A negative transition has the dual
property: its input may be increased while keeping its output fixed.
These are statements about the endpoint relation itself. They need
not say that the entire relation is monotone in both endpoints.

To say that a scalar passage is given by a formula means that its
endpoint membership can be expressed using integer addition, order,
Boolean operations, and quantifiers. This supplies effective
composition and projection without supplying arbitrary transitive
closure.

Presburger sets and relations are as in Definition~\ref{def:presburger}.

The next definition distinguishes the endpoint that may be increased:
the target for a positive passage, and the source for a negative one.

\begin{definition}[One-sided scalar pumping]
\label{def:scalar-one-sided}
Let \(Y=Q_1\times\N\), where \(Q_1\) is finite, and let \(D>0\).
Relations \(P,N\subseteq Y\times Y\) have positive and negative
one-sided pumping, respectively, if
\[
\begin{split}
 ((p,n),(q,m))\in P
   &\ \Longrightarrow\
      ((p,n),(q,m+D))\in P,\\
 ((p,n),(q,m))\in N
   &\ \Longrightarrow\
      ((p,n+D),(q,m))\in N
\end{split}
\]
for all \(p,q\in Q_1\) and all \(n,m\in\N\).
\end{definition}

\begin{theorem}[Scalar saturation]
\label{thm:formal-scalar}
Suppose \(B,P,N\subseteq Y\times Y\) are effectively Presburger
relations, \(B\) is reflexive and transitive, and \(P,N\) satisfy
Definition~\ref{def:scalar-one-sided}. Then, given unary Presburger
sets \(I,J\subseteq Y\), one can decide
\[
 \bigl(I;(B\cup P\cup N)^*\bigr)\cap J\neq\varnothing.
\]
\end{theorem}
\begin{proof}
The proof uses the same two directions as
Theorem~\ref{thm:formal-completion}, with Presburger sets in place of
box sets. We give the formulas to specify precisely which computations
are performed.

Put \(C_-=(N;B)^*\). For any given unary Presburger target \(T\),
iterate
\[
 T_0=T,\qquad
 T_{n+1}=T\cup N^{-1}\bigl(B^{-1}(T_n)\bigr).
 \tag{SB}
\]
Every step is effective Presburger computation. Every added part is
\(D\)-upward by the source-pumping property of \(N\). These added
parts form an increasing sequence, so it stabilizes, with Presburger
equality providing the stopping test. The result is
\(C_-^{-1}(T)\).

For a unary Presburger source \(S\), the set
\(F_1(S)=S;B;C_-;P\) is \(D\)-upward. To compute its basis, use
\[
 F_1(S)\cap H\neq\varnothing
 \quad\Longleftrightarrow\quad
 (S;B)\cap C_-^{-1}\bigl(P^{-1}(H)\bigr)\neq\varnothing
 \tag{SQ}
\]
for a unary Presburger test set \(H\). The right side is effective:
compute the Presburger \(P\)-preimage, use (SB), and test Presburger
intersection. Every principal \(D\)-upward set and its complement is
Presburger, so the basis algorithm of
Lemma~\ref{lem:effective-basis} applies verbatim.
It computes \(F_1(S)\) as a finite union of residue tails.

Now iterate
\[
 S_0=I,\qquad S_{n+1}=I\cup F_1(S_n).
 \tag{SF}
\]
Again the increasing added parts are \(D\)-upward, so the iteration
terminates effectively. Its limit is \(I;A_1\), where
\(A_1=(B;C_-;P)^*\).

Reflexivity and transitivity of \(B\) imply
\((B\cup N)^*=B;(N;B)^*=B;C_-\).
Splitting a word immediately after each \(P\) consequently gives
\[
 (B\cup P\cup N)^*=A_1;B;C_-.
\]
The required answer is the Presburger emptiness test
\[
 (I;A_1;B)\cap C_-^{-1}(J)\neq\varnothing.
\]
The forward iteration computes \(I;A_1\), and the backward iteration
computes \(C_-^{-1}(J)\).
\end{proof}

The conclusion is a set-to-set decision procedure. The proof does not
produce a Presburger formula for the full binary relation
\((B\cup P\cup N)^*\). It also does not assert an effective description
of every forward image or every backward preimage under that full
closure. Only the particular intermediate images and preimages used
in the displayed algorithm have been computed.


\section{Why two bounded counters limit large excursions}\label{sec:pair-rigorous}
% Standalone replacement material. Requires amsthm environments
% definition and lemma, and the usual \N and \Z macros.

\subsection{A numerical measure of distance from an axis}

A configuration can have a very large total even though two counters
are small. Near an axis, it is those two small counters that matter.
We therefore measure their sum. We also apply the same expression to
signed update vectors; in that case it can be negative.

\begin{definition}[Minimum pair sum]\label{def:minimum-pair}
For \(x=(x_1,x_2,x_3)\in\mathbb R^3\), put
\[
 s_2(x)=\min_{1\le i<j\le3}(x_i+x_j),
 \qquad t(x)=x_1+x_2+x_3.
\]
For a finite word of updates \(w=\delta_1\cdots\delta_k\), define
\[
 \Delta(w)=\sum_{j=1}^k\delta_j,\qquad
 t(w)=\sum_{j=1}^k t(\delta_j),\qquad
 a(w)=\sum_{j=1}^k s_2(\delta_j).
\]
The empty word has all three values zero.
\end{definition}

For \(x\in\N^3\), \(s_2(x)\) is the sum of its two smallest entries.
In particular, \(s_2(0,u,v)=\min(u,v)\).
Thus a bound on \(s_2\) bounds two coordinates, while allowing the third
to be arbitrarily large. Notice that \(a(w)\) is a sum of
\emph{step weights}; it is not defined as \(s_2(\Delta(w))\).

The reason for summing the step weights is that they give a lower bound
which can be added along a path.

\begin{lemma}[Adding minimum pair sums]\label{lem:pair-superadditive}
For \(x,y\in\mathbb R^3\),
\[
 s_2(x+y)\ge s_2(x)+s_2(y).
\]
Consequently, if a run with update word \(w\) takes \(x\) to \(y\), then
\[
 s_2(y)-s_2(x)\ge a(w).
\]
\end{lemma}
\begin{proof}
For each fixed pair \(i<j\),
\[
 (x_i+y_i)+(x_j+y_j)
   =(x_i+x_j)+(y_i+y_j)
   \ge s_2(x)+s_2(y).
\]
Take the minimum over the three pairs. Applying this inequality once
per update proves the second assertion.
\end{proof}

\subsection{An exact cycle condition implying pair growth}

We next isolate the finite-graph fact used to obtain pair growth.
The hypothesis is a condition on \emph{all cycles of a specified finite
graph}. It is not a consequence of semilinearity alone, and the graph
must be specified before this lemma is applied.

A directed update graph is a finite directed multigraph \(G=(Q,A)\)
in which each edge \(e\in A\) carries an update
\(\delta(e)\in\Z^3\). A directed walk records only a sequence of control
edges. It need not be enabled at every counter vector. Its update
word, effect, and weights are defined as above.
A simple cycle is a nonempty closed walk whose vertices, other than
its coinciding first and last vertices, are distinct.

The useful situation is that every cycle which gains total also gains
a positive amount of pair weight. A sufficiently small positive
fraction of the total gain can then be charged against pair growth.

\begin{lemma}[Pair growth from cycle weights]\label{lem:cycle-pair-growth}
Let \(G\) be a finite update graph with a nonempty vertex set. Suppose every simple cycle
\(\gamma\) satisfies
\[
 t(\gamma)\ge0,\qquad
 \begin{cases}
 a(\gamma)\ge0,&t(\gamma)=0,\\
 a(\gamma)>0,&t(\gamma)>0.
 \end{cases}
\]
One can compute rational numbers \(\varepsilon>0\) and \(R\ge0\)
such that every enabled run \(x\longrightarrow y\) whose control walk
is in \(G\) satisfies
\begin{equation}
 s_2(y)-s_2(x)
   \ge \varepsilon\bigl(t(y)-t(x)\bigr)-R.
 \label{eq:rigorous-pair-growth}
\end{equation}
\end{lemma}
\begin{proof}
There are finitely many simple cycles. If at least one has positive
total, choose
\[
 0<\varepsilon\le
 \min_{\substack{\gamma\text{ simple cycle}\\t(\gamma)>0}}
       \frac{a(\gamma)}{t(\gamma)}.
\]
Otherwise put \(\varepsilon=1\).
Assign an edge \(e\) the rational weight
\[
 b(e)=s_2(\delta(e))-\varepsilon t(\delta(e)).
\]
Every simple cycle has nonnegative \(b\)-weight. Every closed walk
decomposes into simple cycles, so the graph has no negative-weight
closed walk.

Add a new source with a zero-weight edge to every vertex. Let \(d(q)\)
be the minimum weight of a walk from the new source to \(q\).
These numbers are finite and computable because there is no negative
cycle. For every edge \(e:p\to q\),
\[
 d(q)\le d(p)+b(e).
\]
Hence a walk from \(p\) to \(q\) has \(b\)-weight at least
\(d(q)-d(p)\). With
\[
 R=\max_{q\in Q}d(q)-\min_{q\in Q}d(q)
\]
we obtain \(a(w)\ge\varepsilon t(w)-R\) for every walk \(w\).
Now use Lemma~\ref{lem:pair-superadditive} and
\(t(w)=t(y)-t(x)\).
\end{proof}

This proof is an algorithm: enumerate simple cycles, select
\(\varepsilon\), and compute shortest-path distances in a finite
weighted graph. It does not require knowing whether individual
control cycles are enabled at the endpoints of the query.

\subsection{What a zero pair-weight cycle actually gives}

A cycle of zero pair weight and positive total is excluded by the
strict inequality in Lemma~\ref{lem:cycle-pair-growth}. Symmetry gives
a concrete reason to treat it separately: it yields an actual word
which increases just one counter.

Here symmetry means that every update edge \(p\xrightarrow{\delta}q\)
is accompanied by \(p\xrightarrow{\pi(\delta)}q\) for every permutation
\(\pi\) of the three coordinates. Permutations may be chosen
independently at different steps.

\begin{lemma}[An actual one-coordinate increasing loop]
\label{lem:zero-pair-loop}
In a symmetric update graph, suppose a closed control walk \(w\)
based at \(q\) satisfies \(a(w)=0\) and \(t(w)=T>0\).
There is an effectively constructible closed update word at \(q\)
of effect \(2T e_i\), for each \(i\in\{1,2,3\}\).
There is also an effectively computable \(H\) such that each of these
words is enabled at every \(x\in\N^3\) with \(x_j\ge H\) for all \(j\).
\end{lemma}
\begin{proof}
At every step, put the two smallest update entries in coordinates
1 and 2. The resulting word \(v\) is a closed control walk with
effect \((r,-r,T)\) for some \(r\in\Z\): the first two coordinates
sum to \(a(w)=0\), and the total is \(T\).
Let \(\bar v\) be obtained by swapping coordinates 1 and 2 at every
step of \(v\). Then
\[
 \Delta(v\bar v)=(r,-r,T)+(-r,r,T)=(0,0,2T).
\]
Permuting this complete word gives the other two effects.

For a fixed word \(u\), let \(\Delta(u_{\le k})\) be its effect after
its first \(k\) steps, including the empty prefix. It is enabled from
every vector above the coordinatewise buffer
\[
 B_i(u)=\max_k\bigl\{-\Delta(u_{\le k})_i\bigr\}.
\]
Take \(H\) to be the maximum of all these finitely many buffer
coordinates over the three words.
\end{proof}

Because each resulting word has nonnegative effect, it can be
repeated from the same threshold. The three words can also be
concatenated in any order. Thus, with \(D=2T\), every sufficiently
large \(x\) can reach every \(x+D z\), \(z\in\N^3\), at control \(q\).
This conclusion concerns actual enabled words; it is stronger than
merely finding a coordinate-axis vector in a real cone.

This lemma does not by itself prove that removing the corresponding
controls gives a reachability-preserving construction. That is a
separate requirement of the normalization step.

\subsection{The finite-graph counting lemma}

A large excursion increases a numerical size, but the shorter
connections between excursions may decrease it. The following lemma
states exactly when those decreases can be paid for. It applies to
relations of arbitrary complexity; only their stated numerical bounds
are used.

\begin{definition}[A finite graph of relations with a size function]
\label{def:relation-size-graph}
The data consist of:
\begin{itemize}
\item a finite nonempty set \(Q\) of modes;
\item for each \(q\in Q\), a set \(X_q\) and a function
      \(f_q:X_q\to\mathbb R_{\ge0}\);
\item a finite directed multigraph on \(Q\), with its edges partitioned
      into ordinary edges \(A_0\) and distinguished edges \(A_1\);
\item for every edge \(e:p\to q\), a relation
      \(R_e\subseteq X_p\times X_q\).
\end{itemize}
A realized path is a control walk \(e_1\cdots e_k\) together with
\(x_j\in X_{q_j}\), \(0\le j\le k\), satisfying
\((x_{j-1},x_j)\in R_{e_j}\) for all \(j\).
The number of distinguished edges of this path counts occurrences
of edges in \(A_1\), with repetition.
\end{definition}

In the application, \(X_q\) is a two-counter face, \(f_q\) is a
minimum-pair size, and a distinguished edge is a complete large
excursion. The word \emph{complete} matters: cutting a fixed implementing
word in its middle is handled by an endpoint displacement bound,
not by silently changing this finite graph.

\begin{lemma}[Counting distinguished edges]\label{lem:relation-count}
Take a graph of relations as in Definition~\ref{def:relation-size-graph},
and put \(m=|Q|\). Suppose rational numbers \(P\ge0\), \(\Gamma\),
and weights \(c(e)\), \(e\in A_0\), satisfy:
\begin{enumerate}
\item for each ordinary edge \(e:p\to q\),
\[
 c(e)\ge-P,\qquad
 (x,y)\in R_e\ \Longrightarrow\
       f_q(y)-f_p(x)\ge c(e);
\]
\item every directed cycle consisting only of ordinary edges has
      nonnegative sum of the weights \(c(e)\);
\item for each distinguished edge \(e:p\to q\),
\[
 (x,y)\in R_e\ \Longrightarrow\
       f_q(y)-f_p(x)\ge\Gamma;
\]
\item \(\Gamma-1\ge(m-1)P\).
\end{enumerate}
One can compute \(h:Q\to\mathbb Q\) such that
\(F(q,x)=f_q(x)+h(q)\) is nondecreasing along every realized edge,
and increases by at least one along each distinguished edge.
Consequently, every realized path whose final size is at most \(M\)
has at most
\begin{equation}
 M+\max_{q\in Q}h(q)-\min_{q\in Q}h(q)
 \label{eq:explicit-count}
\end{equation}
distinguished edges.
\end{lemma}
\begin{proof}
Give an ordinary edge weight \(c(e)\), and a distinguished edge
weight \(\Gamma-1\). We first show that this finite weighted graph
has no negative cycle.

An all-ordinary cycle is nonnegative by assumption.
For a cycle containing \(k\ge1\) distinguished edges, split the
ordinary edges into the \(k\) walks between consecutive distinguished
edges. From each ordinary walk delete repeated-vertex closed
subwalks. Each deleted subwalk has nonnegative weight, so the
original walk has weight at least that of the remaining simple path.
The latter has at most \(m-1\) edges and therefore weight at least
\(-(m-1)P\). The original cycle consequently has weight at least
\[
 k\bigl((\Gamma-1)-(m-1)P\bigr)\ge0.
\]

As in Lemma~\ref{lem:cycle-pair-growth}, attach a new zero-weight source
and compute finite shortest-path distances \(d(q)\).
Set \(h(q)=-d(q)\). For an edge \(e:p\to q\) of formal weight \(b(e)\),
\[
 b(e)+h(q)-h(p)
   =b(e)-d(q)+d(p)\ge0.
\]
For an ordinary realized edge, its actual \(f\)-increase is at least
\(b(e)\), so its \(F\)-increase is nonnegative. For a distinguished
edge, its actual \(f\)-increase is at least \(\Gamma=b(e)+1\), so its
\(F\)-increase is at least one.

If a realized path has \(k\) distinguished edges, summing gives
\[
 k\le F(q_{\mathrm{end}},x_{\mathrm{end}})
       -F(q_{\mathrm{start}},x_{\mathrm{start}}).
\]
Use \(f_{q_{\mathrm{end}}}(x_{\mathrm{end}})\le M\) and
\(f_{q_{\mathrm{start}}}(x_{\mathrm{start}})\ge0\) to obtain
\eqref{eq:explicit-count}.
\end{proof}

\subsection{Applying the counting lemma to faces}

On a face, one physical coordinate is stored as a bounded integer
and the other two remain counters. We use a size function whose
change under a two-coordinate translation has an immediate lower
bound.

\begin{definition}[A physical face and its size]\label{def:face-size}
Fix \(B\in\N\). A physical face mode \(q\) specifies a control state
of the original machine, an index \(i_q\in\{1,2,3\}\), and an integer
\(b_q\in\{0,\ldots,B\}\).
Its coordinates are \(z=(z_1,z_2)\in\N^2\), placed in the two physical
positions other than \(i_q\), in increasing order of their indices.
Write \(\iota_q(z)\in\N^3\) for the resulting physical counter vector,
whose coordinate \(i_q\) is \(b_q\). Define
\[
 f_q(z)=b_q+\min(z_1,z_2).
\]
\end{definition}

For every face mode and every \(z\in\N^2\),
\begin{equation}
 0\le f_q(z)-s_2(\iota_q(z))\le B.
 \label{eq:face-size-comparison}
\end{equation}
Indeed the two expressions agree unless both free values are smaller
than \(b_q\); in the remaining case their difference is
\(b_q-\max(z_1,z_2)\).

Suppose an ordinary edge from \(p\) to \(q\) has a fixed translation
of the free pair, \(z\mapsto z+u\), \(u=(u_1,u_2)\in\Z^2\).
Its domain may be restricted by enabling conditions, but every
related pair has that translation. Then
\[
 f_q(z+u)-f_p(z)
   \ge b_q-b_p+\min(u_1,u_2).
\]
This supplies its ordinary weight \(c(e)\).
For a distinguished relation, suppose every related physical pair
\(x=\iota_p(z)\), \(y=\iota_q(z')\) satisfies
\[
 s_2(y)-s_2(x)\ge\varepsilon\bigl(t(y)-t(x)\bigr)-R,
 \qquad t(y)-t(x)\ge C.
\]
By \eqref{eq:face-size-comparison}, its size increase is at least
\[
 \Gamma=\varepsilon C-R-B.
\]
Therefore Lemma~\ref{lem:relation-count} applies as soon as the
ordinary graph satisfies its cycle hypothesis and
\[
 \varepsilon C-R-B-1\ge(m-1)P.
\]
If the final free pair has \(z_1\le A\) or \(z_2\le A\), then
\(f_q(z)\le B+A\); substitute \(M=B+A\) in
\eqref{eq:explicit-count}. The other free coordinate can be arbitrary.

The negative-direction statement uses the reversed relations.
Apply the same lemma to a reversed path: a bound on the physical
pair at the \emph{source} of the original negative path bounds its
number of distinguished edges.

\paragraph{The exact obligation left to the structural construction.}
The preceding lemmas are complete finite-graph arguments. To use them
for a machine, one must separately construct a single finite graph of
relations, prove that it preserves the required runs, and verify the
ordinary-cycle hypothesis. In particular, exhibiting pump words
does not justify deleting their controls; proving a numeric estimate
for one update component does not justify inserting arbitrary other
relations between its paths; and choosing \(C\) after \(m,P,R,B\)
requires those quantities to have been fixed independently of \(C\).
These are hypotheses to verify in the normalization, not conclusions
of the counting lemma.


\section{How the pair bound gives an algorithm for boundary passages}
\label{sec:face-query}

The preceding counting lemma is useful only if its finitely many
patterns can be evaluated. This requires formulas for the individual
relations and for arbitrarily long sequences of ordinary steps.
We now list these requirements explicitly. The result is a theorem
about given finite data; constructing such data from a machine is a
separate task.

\subsection{The exact class of two-counter relation systems}

Think of a directed edge as an instruction to choose \emph{any} pair
from its specified relation. A walk is possible when the intermediate
tuples agree. This semantics makes no reference to an unspecified
notion of a passage or an excursion.

\begin{definition}[Relations specified by a finite graph]
\label{def:relational-graph}
Let \(G\) be a finite directed multigraph with vertices \(Q_G\).
Assign to vertex \(q\) a space \(X_q=\{q\}\times\N^{d(q)}\) and
to every edge \(e:p\to q\) a relation \(R_e\subseteq X_p\times X_q\).
For a walk \(w=e_1\cdots e_k\), its relation is
\[
 R_w=R_{e_1};\cdots;R_{e_k}.
\]
For an empty walk at \(q\), it is the identity on \(X_q\).
For specified start and end vertices, the relation of all accepted
walks is the union of the corresponding \(R_w\)'s.
Equivalently, without start/end restrictions, it is
\((\bigcup_e R_e)^*\).

Finite extra control, such as a sign flag or a position in a fixed
word, is represented by taking a finite product with the vertex set.
Thus restrictions on the sequence of edge types must be given by a
finite automaton or by an explicitly stated relation expression.
\end{definition}

One more property is needed for targets extending to infinity.
Increasing a source need not preserve the same target, but it should
allow some larger target. This weaker property is exactly what makes
preimages of upward sets upward.

\begin{definition}[Weak monotonicity on the grid]
\label{def:weak-monotonicity}
For \(X\) and \(\preceq_D\) as in
Definition~\ref{def:grid-order}, a relation \(R\subseteq X^2\)
is weakly monotone if
\[
 xRy,\quad x\preceq_D x'
 \quad\Longrightarrow\quad
 \exists y'\in X:\ x'Ry'\ \text{and }y\preceq_D y'.
 \tag{WM}
\]
Labels on each comparison agree by the definition of \(\preceq_D\).
\end{definition}

Indeed, if \(T\) is upward, \(xRy\in T\), and \(x\preceq_D x'\),
property (WM) provides \(x'Ry'\) with \(y'\in T\).
It follows that \(R^{-1}(T)\) is upward. Merely knowing that \(R\)
has a Presburger formula does not give (WM).

\begin{theorem}[Two-dimensional box queries]
\label{thm:face-query}
Suppose a finite graph of relations has the following given data and
properties.
\begin{enumerate}[label=\textnormal{(F\arabic*)},leftmargin=*]
\item Each vertex is a physical face mode from
Definition~\ref{def:face-size}, with stored value at most \(B\).
Every edge relation \(R_e\) is effectively Presburger.
\item The edges are partitioned into ordinary and distinguished
edges. Put \(O=\bigcup_{e\in A_0}R_e\) and
\(F=\bigcup_{e\in A_1}R_e\). A Presburger formula for \(O^*\)
is supplied.
\item The hypotheses of Lemma~\ref{lem:relation-count} hold for the
face size \(f_q(z)=b_q+\min(z_1,z_2)\), with all its constants supplied.
\item Every edge relation is weakly monotone on a fixed \(D\)-grid.
\end{enumerate}
Let \(\Phi=(O\cup F)^*\). Given a Presburger source set \(S\) and a
grid-box target set \(T\), one can decide
\((S;\Phi)\cap T\ne\varnothing\).
\end{theorem}
\begin{proof}
It suffices to consider one target box.

\emph{Case 1: one target coordinate is bounded.}
Suppose \(z_i\le A\) on the box. Then \(f_q(z)\le B+A\) there.
Lemma~\ref{lem:relation-count} gives an integer bound \(L\) on the
number of distinguished edges in every relevant walk. Let
\[
 H_L=\bigcup_{k=0}^{L} O^*;(F;O^*)^k.
 \tag{FP}
\]
Every factor in this finite expression has a Presburger formula, so
\(H_L\) is effectively Presburger. On endpoints in \(S\times T\),
the relations \(H_L\) and \(\Phi\) agree. A Presburger emptiness test
answers the question.

\emph{Case 2: neither target coordinate is bounded.}
The box \(T\) is upward. Put \(R=O\cup F\) and iterate
\[
 T_0=T,\qquad T_{n+1}=T\cup R^{-1}(T_n).
 \tag{UP}
\]
Each term is effectively Presburger. By (WM), every term is upward.
The sequence is increasing, hence stabilizes by Dickson's lemma;
Presburger equality detects the first fixed point.
Induction gives \(T_n=(R^{\le n})^{-1}(T)\), so the fixed point is
\(\Phi^{-1}(T)\). Intersect it with \(S\).
\end{proof}

This is the precise reason that boxes are useful. A box in \(\N^2\)
either has an upper bound on a coordinate or is upward within its
residue class. A general Presburger target need not satisfy either
alternative: the diagonal \(\{(n,n):n\in\N\}\) does not.

The proof also gives two stronger, but restricted, outputs.
In Case 1 it gives the full relation restricted to that target box.
In Case 2 it gives the predecessor set of that particular target.
It does not give the full relation \(\Phi\).

\subsection{Why one-dimensional endpoints are particularly useful}

An axis-near endpoint has two bounded \emph{physical} coordinates.
On a face record one of them is already stored. Consequently one of
the two free coordinates is bounded as well. The bound is fixed by the
representation, not by the remaining unbounded value.

\begin{corollary}[A uniform finite pattern bound at an axis]
\label{cor:axis-patterns}
Under the hypotheses of Theorem~\ref{thm:face-query}, fix \(W\in\N\).
The restriction of \(\Phi\) to targets satisfying
\(\min(z_1,z_2)\le W\) is effectively Presburger.
In particular its relation between any fixed finite family of scalar
axis encodings is effectively Presburger.
\end{corollary}
\begin{proof}
The target size is at most \(B+W\), so the same integer \(L\) works
for all target values. Use (FP) and impose the target condition.
Restricting an endpoint to an axis encoding is a Presburger
substitution of fixed coordinates.
\end{proof}

For negative passages, apply the same reasoning to the graph with
every relation reversed. The bounded physical pair is then required
at the \emph{source} of the original passage.
This direction matters when choosing which endpoint representation to
use; an arbitrary scalar summary of a three-counter vector need not
bound any physical pair.

\subsection{Where the formula for ordinary paths can come from}

A finite list of words is not yet a symmetric two-counter machine:
adding the coordinate-swapped version of each update also permits
mixed choices of swaps at different positions. The exact construction
must account for those choices.

\begin{definition}[A symmetric word implementation]
\label{def:word-implementation}
For a fixed nonempty two-counter update word
\(w=\eta_1\cdots\eta_k\), \(k\ge1\), let
\[
 \operatorname{Orient}(w)
   =\{\pi_1(\eta_1)\cdots\pi_k(\eta_k):
                         \pi_j\in S_2\}.
\]
Its private-chain implementation introduces distinct intermediate
controls for positions \(1,\ldots,k-1\). At position \(j\), it places
both coordinate permutations of \(\eta_j\) between the prescribed
controls. No other edge touches an intermediate control.
\end{definition}

A prescribed identity relation can instead be implemented by a single
zero-vector update edge. Every run between the ends of a nonempty-word
chain reads exactly one word in
\(\operatorname{Orient}(w)\), and every enabled such word is a run of
the chain. Taking disjoint chains for a finite family of words gives
an exact symmetric two-VASS implementation of their concatenations.
Assumption \textbf{S2} therefore computes its reachability relation
at the designated endpoints. If a proposed ordinary relation has
extra tests, marker changes, or coordinate recodings, an exact
implementation of those operations must also be supplied; the
private-chain argument alone does not supply them.



\section{What the auxiliary records mean physically}
\label{sec:records}

The numerical algorithms operate on labelled tuples. To apply them to
a machine, every tuple that represents a physical configuration must
have an explicit interpretation. A label can store finitely many
bounded counter values. It cannot store an unbounded quantity unless
that quantity appears as one of the tuple's coordinates.

\subsection{Bounded coordinates and one-counter records}

A zero coordinate is just one special case of a bounded coordinate.
Fixing one bounded coordinate leaves two free coordinates; fixing two
leaves one. The bounds must be fixed before analyzing an arbitrarily
long run.

\begin{definition}[Thin regions and axis records]
\label{def:axis-record}
For \(H,W\in\N\), put
\[
\begin{split}
 \operatorname{Thin}_H
   &=\{(q,x):\min_i x_i\le H\},\\
 \operatorname{Axis}_W
   &=\{(q,x):\exists i<j,\ x_i\le W,\ x_j\le W\},\\
 \operatorname{High}_H
   &=\{(q,x):x_i\ge H\text{ for all }i\}.
\end{split}
\]
A physical axis record has label
\(a=(q,i,j,b,c)\), with \(i<j\) and \(0\le b,c\le W\), and one
coordinate \(n\in\N\). It denotes \((q,x)\) with
\[
 x_i=b,\qquad x_j=c,\qquad x_k=n
 \quad\text{for }\{k\}=\{1,2,3\}\setminus\{i,j\}.
\]
Its physical total is \(n+b+c\).
A physical face record is the record of
Definition~\ref{def:face-size}; the bounded coordinate chosen there
is also called its stored coordinate.
\end{definition}

All labels here range over finite sets. An axis record has an actual
bounded physical pair, independently of \(n\). This is the property
used in Corollary~\ref{cor:axis-patterns}.

Here is a small but useful observation about changing which coordinate
is stored. Let
\[
 J=\max\{|\delta_i|:(p,\delta,q)\in\Delta,\ 1\le i\le3\},
\]
with \(J=0\) if there are no update transitions.
If an update changes \(x\) to \(y\), \(x_i\le H\), and \(y_j\le H\)
for distinct \(i,j\), then
\[
 x_j\le H+J,\qquad y_i\le H+J.
\]
Thus both endpoints lie in \(\operatorname{Axis}_{H+J}\).
A zero test of a free coordinate also reaches an axis record, because
the stored coordinate is bounded and the tested one is zero.
These facts explain the intended separation into fixed-stored-coordinate
intervals and scalar records. They do not alone construct a simulator
for all changes of representation.

\subsection{What an exact two-counter presentation must provide}

The phrase that a part of the construction is handled by S2 needs
more than the observation that two numbers are visible at its ends.
It needs an actual finite symmetric two-counter machine and an
endpoint encoding for which reachability is exactly the intended
relation.

\begin{definition}[Exact two-counter presentation]
\label{def:exact-presentation}
Let \(T\subseteq X^2\) be a relation on a finite union of labelled
counter spaces. An exact symmetric two-counter presentation of \(T\)
consists of a specified symmetric two-counter machine \(N\) and a
specified injective encoding
\[
 \eta:X\longrightarrow \operatorname{Conf}(N)
\]
whose graph is effectively Presburger, such that
\[
 uTv\quad\Longleftrightarrow\quad
 \eta(u)\longrightarrow_N^*\eta(v)
 \qquad(u,v\in X).
 \tag{EX}
\]
If only a designated subset of \(X\) is encoded, both the domain and
the corresponding restriction of (EX) must be specified.
\end{definition}

Under S2, (EX) gives a formula for \(T\) by substituting the endpoint
encoding into the computed reachability formula of \(N\).
The forward implication in (EX) prevents losing intended runs.
The reverse implication prevents extra permutations of updates in
the simulator from creating unintended runs.
Neither implication is supplied merely by naming \(N\) a simulator.

\subsection{A scalar total is a different kind of record}

Sometimes different three-counter vectors with the same total can
reach one another at a particular control state. Then one integer
can stand for an entire class of vectors. This is useful, but it is
not the same representation as an axis record: the class may contain
vectors very far from every axis.

\begin{definition}[Transfer records]
\label{def:transfer-record}
Fix a control \(q\), a threshold \(H\), a modulus \(m>0\), and
\(r\in\{0,\ldots,m-1\}^3\). For \(T\in\N\), define
\[
 \mathcal F_{q,H,m,r}(T)
 =\{(q,x):x_i\ge H,\ x_i\equiv r_i\pmod m\ (1\le i\le3),
                        \ \textstyle\sum_i x_i=T\}.
\]
The parameters give valid transfer records if every nonempty such
set satisfies
\[
 c,d\in\mathcal F_{q,H,m,r}(T)
       \quad\Longrightarrow\quad c\longrightarrow_M^*d.
 \tag{TR}
\]
The record \((q,H,m,r,T)\) denotes this set of configurations.
Only values of \(T\) for which the set is nonempty are legal.
The control \(q\) for which (TR) is proved is called its transfer
anchor. No property is asserted at another control.
\end{definition}

For each \(i\), let \(b_i\) be the least integer at least \(H\) with
residue \(r_i\). The set in the definition is nonempty exactly when
\[
 T\ge b_1+b_2+b_3,\qquad
 T\equiv b_1+b_2+b_3\pmod m.
\]
When it is nonempty, a canonical vector in it is
\((b_1,b_2,T-b_1-b_2)\).
Property (TR) justifies passage between this vector and every other
vector in the set.

It does \emph{not} justify transferring a pair-size bound from the
canonical vector to another representative. For example, when
\(m=1\), a fixed \(H\) and large \(T\) allow both
\((H,H,T-2H)\) and, for \(T\) divisible by three,
\((T/3,T/3,T/3)\). Their minimum pair sums are \(2H\) and \(2T/3\).
Hence a bounded pair at the canonical representation can conceal an
arbitrarily large physical pair. This is why a transfer record cannot
be used as an axis endpoint in the counting lemma without a separate
argument.

The following observation explains exactly how a set-valued
representation can still be used soundly in compositions.

\begin{lemma}[Composing through mutually reachable representatives]
\label{lem:fibre-composition}
Suppose each legal record \(u\) denotes a nonempty set
\(\rho(u)\subseteq\operatorname{Conf}(M)\), and every two members of
\(\rho(u)\) are mutually reachable. Suppose every pair \(uRv\) has
some witnesses \(c\in\rho(u)\), \(d\in\rho(v)\) with
\(c\longrightarrow_M^*d\).
Then a finite sequence of \(R\)-steps between records can be realized
from any representative of its first record to any representative of
its last record.
\end{lemma}
\begin{proof}
Choose the witnessing machine path for each step. At each join the
previous witness ends and the next witness starts in the same set
\(\rho(u)\). Property (TR), or its stated analogue for that record,
connects those two configurations. The same property connects the
prescribed outer representatives to the first and last witnesses.
\end{proof}

\subsection{Relations used to enter or leave a representation}

An adapter is not a new kind of computational power. It is a binary
relation with specified source and target representations. For example,
the change from a face record to an axis record is the relation saying
that both records denote the same physical vector.

\begin{definition}[Specified run relations and representation changes]
\label{def:run-relations}
Let \(X_p,X_q\) be two record spaces with interpretation sets
\(\rho_p(x),\rho_q(y)\).
For a specified family \(\mathcal W\) of allowed machine runs, define
\[
 R^{\mathcal W}_{pq}
 =\{(x,y):\exists c\in\rho_p(x),\ d\in\rho_q(y),\
             \text{a run in }\mathcal W\text{ goes from }c\text{ to }d\}.
\]
A relation of this form is called a run relation.
The family \(\mathcal W\) must be specified, for example by an
update graph with fixed entry and exit words, or by a concrete
two-counter presentation.
A change of representation without a machine step is the relation
\[
 C_{pq}=\{(x,y):\rho_p(x)\cap\rho_q(y)\ne\varnothing\}.
\]
\end{definition}

In a finite graph of relations, an adapter is simply a designated
edge of one of these types, possibly restricted by an explicit
Presburger endpoint condition. Its formula, or an algorithm for the
queries required of it, is part of the construction. The word adapter
does not imply either one.

\subsection{Private operations: exactly when they are sound}

A Cartesian relation forgets how the particular source and target
were paired: every permitted source can reach every permitted target.
It can consequently be split through one auxiliary point. This is
safe only when the entire product, not just its two projections, is
known to be realizable.

\begin{definition}[Cartesian factorization]
\label{def:cartesian}
For \(E\subseteq X\), \(F\subseteq Y\), and a fresh one-point space
\(\{\ast\}\), put
\[
 C_{\mathrm{in}}=\{(x,\ast):x\in E\},\qquad
 C_{\mathrm{out}}=\{(\ast,y):y\in F\}.
\]
Then \(C_{\mathrm{in}};C_{\mathrm{out}}=E\times F\).
The point \(\ast\) is called private if the finite graph has only the
prescribed entrances and exits at it and it is not an original
source or target configuration.
\end{definition}

If \(E\times F\) is contained in the intended machine reachability
relation, the complete block is sound. In contrast, projecting a
relation \(R\) onto its domain and range gives only
\(R\subseteq\operatorname{dom}(R)\times\operatorname{ran}(R)\);
the reverse inclusion is false even for the identity relation.

A second private operation inserts a formal pump which can be
absorbed into the preceding or following relation. It need not be a
physical loop at the intermediate record.

\begin{lemma}[Absorbing a private pump]
\label{lem:private-pump}
Let \(R\subseteq X\times Y\) and \(U\subseteq Y^2\) satisfy
\(R;U=R\). Replace an \(R\)-edge by an \(R\)-edge to a fresh copy of
\(Y\), followed by a \(U\)-edge back to the original target copy.
If these are the only entrance and exit of the fresh copy, and
original endpoints are outside it, the replacement preserves every
relation between original endpoints.
Dually, \(V;R=R\) permits a mandatory private \(V\)-edge before \(R\).
\end{lemma}
\begin{proof}
Every visit to the fresh copy occurs inside the complete consecutive
block \(R;U\). Substitute the equality \(R;U=R\) in a composed word.
The reverse substitution proves completeness. The dual argument is
identical.
\end{proof}

For instance, let \(Y=Q_1\times\N\) and let \(U_D\) add any
nonnegative multiple of \(D\) without changing its label.
If \(xR(q,n)\) implies \(xR(q,n+D)\), iteration and the zero-increment
case give \(R;U_D=R\). The identity is a condition to prove for the
particular \(R\); allowing another entrance to the private copy would
invalidate this justification.



\section{One fixed relation between the independent events}
\label{sec:fixed-K}

We can now define the relation previously called a connector without
relying on that name. There are external record spaces, where an
independent event may occur, and internal scalar record spaces.
A passage either avoids the internal scalar spaces or has a first
and a last scalar record. The part between them may contain
arbitrarily many scalar steps.

\begin{definition}[Data defining the fixed relation \(K\)]
\label{def:K-data}
Let \(X\) be a finite union of counter spaces of dimensions at most
three, and let \(A=Q_A\times\N\), with \(Q_A\) finite.
Treat \(X\) and \(A\) as disjoint tagged sets, even if some of their
records have the same physical interpretation.
Specify the following relations:
\[
\begin{array}{rcll}
 Z_0&\subseteq&X\times X
   &\text{passages with no selected internal scalar record},\\
 R_{\mathrm{in}}&\subseteq&X\times A
   &\text{passages to the first such record},\\
 R_{\mathrm{out}}&\subseteq&A\times X
   &\text{passages from the last such record},\\
 B_0,\ E&\subseteq&A\times A
   &\text{two specified kinds of scalar passages}.
\end{array}
\]
Their mathematical descriptions are part of the data; the explanations
in the right column are intended uses, not their definitions.
Define
\begin{equation}
 K=Z_0\ \cup\
 R_{\mathrm{in}};(B_0\cup E)^*;R_{\mathrm{out}}.
 \label{eq:defined-K}
\end{equation}
\end{definition}

For example, a description of \(R_{\mathrm{in}}\) may be a finite
graph of relations with specified initial and final vertices, as in
Definition~\ref{def:relational-graph}. A permitted walk in that graph
is then the exact meaning of a passage. Saying instead that it is a
normalized passage, without specifying the graph or an equivalent
relation expression, would leave \(R_{\mathrm{in}}\) undefined.

If an external endpoint is already represented by a scalar record,
the corresponding \(R_{\mathrm{in}}\) or \(R_{\mathrm{out}}\) includes
the explicitly specified identity change between its two tagged
copies. This includes empty prefixes and suffixes.
An internal visit is a visit in this \emph{specified description}.
It does not mean that every physical run realizing an edge must be
split whenever it happens to come near an axis.

\subsection{The scalar property needed in addition to semilinearity}

One scalar passage may be represented by finitely many patterns, so
S2 and S3 can give its Presburger formula. To compose arbitrarily many
such passages, the proof needs an additional property: a sufficiently
large gain permits a further increase of the target while fixing the
source. Large loss has the reversed property.

\begin{definition}[Large-gain and large-loss pumping]
\label{def:large-pumping}
A scalar relation \(E\subseteq(Q_A\times\N)^2\) has effective
large-gain and large-loss pumping if one can compute integers
\(D>0\), \(C\ge0\) such that, for every \(p,q,n,m\),
\[
\begin{split}
 (p,n)E(q,m),\quad m-n>C
   &\ \Longrightarrow\ (p,n)E(q,m+D),\\
 (p,n)E(q,m),\quad n-m>C
   &\ \Longrightarrow\ (p,n+D)E(q,m).
\end{split}
\tag{LP}
\]
For these constants, put
\[
\begin{split}
 E_+&=\{((p,n),(q,m))\in E:m-n>C\},\\
 E_-&=\{((p,n),(q,m))\in E:n-m>C\},\\
 E_0&=\{((p,n),(q,m))\in E:|m-n|\le C\}.
\end{split}
\]
\end{definition}

If \(E\) is Presburger, all three parts are effectively Presburger.
Condition (LP) makes \(E_+\) and \(E_-\) satisfy
Definition~\ref{def:scalar-one-sided}: increasing the indicated
endpoint also preserves the strict gain or loss condition.
The example \(m=2n\) in Section~\ref{sec:model} shows why (LP)
requires a separate proof for the actual run families.

\subsection{Why bounded scalar differences can be added using S2}

In one dimension, a bounded difference means one of finitely many
fixed shifts. The only remaining choice is a unary condition on the
starting value. Such conditions can be checked at an actual two-counter
axis, where the unused counter is zero.

\begin{lemma}[Adding Presburger-guarded bounded shifts]
\label{lem:band-compiler}
Let \(N_0\) be a specified symmetric two-counter machine with distinct
designated controls \(p_a\), \(a\in Q_A\). Interpret
\((a,n)\) as \((p_a,n,0)\), and let \(B_0\) be the restriction of
\(\Reach(N_0)\) to these records.
Let \(E_0\subseteq(Q_A\times\N)^2\) be an effectively Presburger
relation satisfying \(|m-n|\le C\) for a fixed integer \(C\).
One can construct a symmetric two-counter machine whose reachability
on the encoded scalar records is exactly
\[
 B=(B_0\cup E_0)^*.
 \tag{BC}
\]
Consequently \(B\) is effectively Presburger under S2.
\end{lemma}
\begin{proof}
For \(a,b\in Q_A\) and \(-C\le d\le C\), compute the unary guard
\[
 G_{a,b,d}=\{n\in\N:n+d\ge0,\
                     ((a,n),(b,n+d))\in E_0\}.
\]
Every unary Presburger set can effectively be written as finitely many
exceptional values and finitely many tails
\(\{n:n\ge h,\ n\equiv r\pmod M\}\), using one common modulus \(M\).

Refine each control of \(N_0\) by the value of the encoded total
\(x+y\pmod M\). An update \((u,v)\) changes this finite-control value
by \(u+v\pmod M\). Its coordinate swap has the same change, so the
refinement preserves symmetry. At a scalar port \((n,0)\), this
stored value gives \(n\pmod M\). The encoding uses the unique
consistent residue copy of \(p_a\); it is Presburger-definable.

To test \(n\ge h\) without changing \(n\), use \(h\) unit decrements
followed by \(h\) unit increments, with the second counter tested
zero before and after the gadget. Include both coordinate
permutations of every update. During the decrements the second
counter is zero, so a decrement on it is disabled. During the
increments any increase of the second counter would survive until
the final zero test, so every successful run increments only the
first. The gadget therefore exists exactly when \(n\ge h\), and
restores \((n,0)\). To test \(n=h\), additionally test the first
counter zero between the decrements and increments.

Residue conditions are checked by finite control. For each permitted
guard branch, follow its test by a fixed shift \(d\): use \(d\) unit
increments when \(d\ge0\), or \(-d\) unit decrements otherwise,
again testing the second counter zero at the two ends.
The same argument forces the intended orientation. Every completed
gadget takes \((p_a,n,0)\) to \((p_b,n+d,0)\) for exactly the pairs
in that branch.

The total-residue refinement is maintained on every new gadget
update as well as on the original transitions.
All intermediate controls are fresh and have only their prescribed
successors. Thus a path between designated ports alternates paths
of \(N_0\) and complete new gadgets. Its relation is (BC), in both
directions. Apply S2 to the resulting finite symmetric machine and
restrict to the consistently encoded ports.
\end{proof}

This lemma has a specific encoding hypothesis: the scalar ports are
actual axes of the supplied two-counter machine.
A state-dependent scalar convention \(n=x+c_a\) can be handled by
the corresponding fixed substitutions; a step \(m=n+d\) then shifts
the active counter by \(d+c_a-c_b\).
An arbitrary binary Presburger relation is not covered.

There is also a basic dimensional distinction. On a two-dimensional
face, a bound on the \emph{total} difference does not bound individual
coordinate changes: all transformations
\((u,v)\mapsto(u+k,v-k)\), when nonnegative, have total difference
zero. Therefore the bounded-band argument above does not establish
semilinearity of the closure of the original two-dimensional
zero-gain segment relation.

\subsection[The exact algorithm for K]{The exact algorithm for \(K\)}

To evaluate the first and last part of
\eqref{eq:defined-K}, we need formulas for the sets of scalar
values they produce or accept. A yes/no answer for each scalar value
does not by itself compute such a formula. This requirement is
therefore part of the statement.

\begin{theorem}[A sufficient algorithmic specification for \(K\)]
\label{thm:K-algorithm}
Take the fixed relations of Definition~\ref{def:K-data}. Suppose:
\begin{enumerate}[label=\textnormal{(K\arabic*)},leftmargin=*]
\item \(\Hit_{Z_0}(S,T)\) is decidable for all grid-box sets
      \(S,T\subseteq X\);
\item the unary Presburger sets
\[
 I(S)=S;R_{\mathrm{in}},\qquad
 J(T)=R_{\mathrm{out}}^{-1}(T)
\]
      are effectively computable for every such \(S,T\);
\item \(E\) is effectively Presburger and has effective pumping (LP);
\item an effective Presburger formula is supplied for
      \(B=(B_0\cup E_0)^*\).
\end{enumerate}
Then \(K\) has a box-intersection algorithm.
\end{theorem}
\begin{proof}
By the definition of \(K\), its query is exactly
\begin{equation}
 \Hit_K(S,T)
 \quad\Longleftrightarrow\quad
 \Hit_{Z_0}(S,T)\ \lor\
 \bigl(I(S);(B_0\cup E)^*\bigr)\cap J(T)\ne\varnothing.
 \label{eq:exact-K-query}
\end{equation}
Since \(E=E_0\cup E_+\cup E_-\),
\[
 (B_0\cup E)^*=(B\cup E_+\cup E_-)^*.
\]
The relation \(B\) is reflexive and transitive by its definition.
Apply Theorem~\ref{thm:formal-scalar} to \(B,E_+,E_-\), using
the two computed scalar endpoint sets. Together with (K1), this
decides \eqref{eq:exact-K-query}.
\end{proof}

Theorem~\ref{thm:face-query} explains how (K1) and (K2) can be obtained
for specified two-dimensional passage graphs. At a fixed axis
endpoint, retain its free scalar variable in the finite-pattern
formula. For an upward target, compute the predecessor formula and
restrict it to that scalar endpoint. Negative passages use the
reversed calculation. Additional full-dimensional endpoints and
transfer records require their own specified relations and query
proofs; the theorem does not silently cover them.

All descriptions here are fixed before a query. A numerical bound
on the length of the finite expression (FP) may depend on the query
box, but the relation being evaluated is always the same \(K\).



\section{The geometric idea behind the proposed construction}
\label{sec:geometry}

The algorithms above have exact hypotheses. The intended geometric
argument tries to produce relations satisfying them by separating
update-only behavior according to its cycles. This section defines
that geometry and explains its intended role without identifying a
real-cone calculation with an enabled machine run.

\subsection{Cycle effects and their signs}

A long control walk consists of a bounded simple path and repeated
cycles. Nonnegative combinations of cycle effects describe the
directions in which large displacements can occur, ignoring for the
moment enabling constraints and integrality.

\begin{definition}[Cycle cone and signed components]
\label{def:cycle-cone}
Let \(G\) be a finite update graph. A strongly connected component
is a maximal nonempty vertex set in which each vertex is reachable
from each other by a directed control walk.
For such a component \(S\), let \(\mathcal C_S\) be its finite set
of simple directed cycles, with their update labels. Define
\[
 C_S=\left\{\sum_{\gamma\in\mathcal C_S}
          \lambda_\gamma\Delta(\gamma):
                         \lambda_\gamma\in\R_{\ge0}\right\}.
\]
Here a convex cone is a set closed under addition and multiplication
by nonnegative real scalars. A cone \(C\) is pointed if
\(C\cap(-C)=\{0\}\).
Call \(S\) positive if
\(\tot v>0\) for every nonzero \(v\in C_S\), and negative if
\(\tot v<0\) for every such \(v\), where \(\tot v=\sum_i v_i\).
Call a positive component productive if its cone contains a vector
with a negative coordinate. A positive component is called orthant
if \(C_S\subseteq\R_{\ge0}^3\).
\end{definition}

In a symmetric update graph, \(C_S\) is invariant under coordinate
permutations. For a nonzero permutation-invariant pointed cone,
all nonzero vectors have the same strict total sign.
To prove this, a nonzero vector \(v\) of total zero would contradict
pointedness: the sum of all its permutations is zero, so
\(-v\) is a nonnegative sum of the other permutations.
Vectors of both positive and negative total also contradict
pointedness, because averaging their permutations gives positive and
negative multiples of \((1,1,1)\).
This proves the sign assertion as a statement about cones.

The productive case has a direction which increases both free
coordinates on any chosen face. The stored coordinate can then be
left unchanged.

\begin{lemma}[A direction inside a productive cone]
\label{lem:face-interior}
Let \(C\subseteq\R^3\) be a permutation-invariant convex cone such that
every nonzero vector of \(C\) has positive total.
If \(C\) contains a vector with a negative coordinate, then each
\[
 d_1=(0,1,1),\qquad d_2=(1,0,1),\qquad d_3=(1,1,0)
\]
belongs to the interior of \(C\) in \(\R^3\).
Here interior means that some open Euclidean ball about the vector
is contained in the cone.
\end{lemma}
\begin{proof}
Permute a negative entry into the first coordinate and average the
other two entries by their swap. This gives
\(w_1=(-a,b,b)\in C\), where \(a>0\) and \(2b>a\).
Its permutations \(w_2=(b,-a,b)\), \(w_3=(b,b,-a)\) are also in \(C\).
They are linearly independent: their determinant is
\((2b-a)(a+b)^2>0\). Moreover
\[
 \frac{2b}{a}w_1+w_2+w_3
    =\frac{(2b-a)(a+b)}{a}\,(0,1,1).
\]
All coefficients are positive. In the basis \(w_1,w_2,w_3\), every
sufficiently small perturbation of this vector still has positive
coordinates, so it remains in their nonnegative span and hence in
\(C\). Positive scaling and permutation prove the claim.
\end{proof}

This explains the proposed source-to-target compensation in (WM):
after a bounded change of displacement, moving sufficiently far in
an interior direction stays in the cone. What still needs a
machine argument is that the new displacement belongs to the
appropriate integer class and has an \emph{enabled run with the
required controls and restrictions}. A direction in \(C_S\) is a
real combination of cycle effects, not itself a machine transition.

In an orthant component, a different elementary argument is
available. For any coordinate \(i\), the \(i\)-th update weights
have no negative cycle. A finite control potential therefore gives
a constant \(c_i\) such that every subpath has \(i\)-th effect at
least \(-c_i\). If a path ends at \(y\), every earlier vector \(x\)
on it satisfies \(x_i\le y_i+c_i\).
Thus a bounded endpoint coordinate stays bounded throughout that
path. This is the reason orthant pieces can sometimes be assigned
to a bounded-coordinate presentation.

\subsection{Which update-only relations S3 actually describes}

A relation must specify both its graph and any additional conditions
on runs. In particular, one cannot invoke S3 for a relation described
only as all paths avoiding a vaguely specified event.

\begin{definition}[A specified update-only relation with a high visit]
\label{def:specified-update}
Let \(G\) be a specified symmetric subgraph of the update graph,
with inherited counter semantics. For \(p,q\) in \(G\), set
\[
 R^G_{p,q}=\{(x,y):(p,x)\longrightarrow_G^*(q,y)\}.
\]
For a specified control \(r\) and threshold \(H\), set
\[
 R^{G,r,H}_{p,q}
 =\{(x,y):\exists z\in\N^3,\
       z_i\ge H\ (1\le i\le3),\
       xR^G_{p,r}z,\ zR^G_{r,q}y\}.
\]
Fixed entry and exit words, and explicit Presburger endpoint
conditions such as \(\tot y-\tot x>\Co\), are added by finite
relation composition and intersection.
\end{definition}

Assumption S3 computes all these relations. Removing a specified
set of controls, or taking copies of controls while retaining all
coordinate permutations of each update, still gives a symmetric
VASS to which S3 applies. An arbitrary restriction on the numerical
configurations visited by a path is not automatically of this form.

In the proposed construction, a large positive relation should be
one of these specified relations, with a gain condition and a
specified residual graph satisfying Lemma~\ref{lem:cycle-pair-growth}.
The adjective residual means only that the graph has been obtained
by a specified finite removal or control-copy operation; it does
not supply a preservation theorem for the original machine.

\subsection{What changing the order of signs would have to prove}

A positive relation followed by a negative one can sometimes be
connected by making an intermediate vector sufficiently large.
A negative relation followed by a positive one can sometimes be
lowered in the middle until it reaches an axis.
These are two different claims, and each needs an endpoint relation
statement.

\begin{definition}[The two proposed replacement properties]
\label{def:replacements}
Let \(X\) be the finite disjoint union of the physical face-record
spaces under discussion, and define their physical decoding by
\(\rho(q,z)=(\operatorname{control}(q),\iota_q(z))\), with
\(\iota_q\) from Definition~\ref{def:face-size}.
Suppose \(P,N,O\subseteq X^2\) are specified run relations, with
\(O^*\) defined by a specified ordinary graph.
A \emph{Cartesian replacement} for the order \(P,N\) is a finite
list of specified sets \(E_\ell,F_\ell\) such that
\[
 P;O^*;N\ \subseteq\
       \bigcup_\ell E_\ell\times F_\ell
       \ \subseteq\ \Reach_{\rho}(M),
 \tag{PK}
\]
where \(x\Reach_\rho(M)y\) means that the physical configurations
denoted by \(x,y\) are connected by an \(M\)-run.

An \emph{axis replacement} for the order \(N,P\) is a finite list of
specified relations \(L_\ell\subseteq X\times A_\ell\) and
\(R_\ell\subseteq A_\ell\times X\), where \(A_\ell\) consists of
physical axis records, such that
\[
 N;O^*;P\ \subseteq\
       \bigcup_\ell L_\ell;R_\ell
       \ \subseteq\ \Reach_\rho(M).
 \tag{VL}
\]
All sets and relations in the replacement must belong to the
specified effective family of the construction.
\end{definition}

These definitions state what the formerly informal peak and valley
arguments need to deliver. The second inclusion in each display is
soundness; the first ensures that every original occurrence can be
replaced. If one needs equality with the original local relation,
that is a stronger assertion and must be proved separately.
For the overall decision procedure, a sound enlargement can suffice.

The intuitive valley operation is to subtract the same amount from
the two free coordinates of a fixed-coordinate middle segment.
The stored coordinate and its zero tests are then unchanged.
One stops before a free counter becomes negative; an axis record is
obtained. The neighboring signed update pieces must be reconstructed
to the altered endpoints. It is that last sentence which needs a
realization proof in the \emph{same specified update graphs}; the
cone calculation alone does not establish it.

A finite decreasing measure is useful once these replacements have
been established. It does not establish their existence.

\begin{lemma}[The role of the occurrence-label measure]
\label{lem:normal-measure}
Consider a finite word of signed occurrences, each bearing a distinct
original label, with some positions designated as cuts. Suppose
available rewrites have the following properties:
\begin{enumerate}[leftmargin=*]
\item a rewrite of a consecutive positive/negative pair removes
      those two labels and creates no new labels;
\item a rewrite of a consecutive negative/positive pair inserts a
      cut between the two labels, preserving the labels and signs;
\item no rewrite removes cuts, changes the order or signs of retained
      labels, or inserts a new signed occurrence between them.
\end{enumerate}
Then repeated rewrites of unlike adjacent signs terminate.
\end{lemma}
\begin{proof}
Order lexicographically the pair consisting of the number of
retained labels and the number of unlike-sign adjacencies within
uncut intervals. The first rewrite lowers the first coordinate.
The second preserves it and removes at least the selected adjacency,
without creating a new one. The pair lies in \(\N^2\), whose
lexicographic order has no infinite strictly decreasing sequence.
\end{proof}

To apply this lemma to runs, one must prove that each numerical
replacement has exactly these effects on the chosen description.
For example, a reconstructed update-only relation may have a
realizing run with extra physical boundary visits. If it is retained
as one edge of the fixed relation graph, those visits are not new
letters. If it is expanded again, the measure proof above no longer
applies without further analysis.



\section{The remaining construction, stated without informal placeholders}
\label{sec:construction}

The previous sections prove algorithms for precisely described
relations. To decide reachability of the original machine, those
relations must now be produced from its finite transition table.
This is a mathematical construction theorem, not a definition of
another name for the original reachability relation.

\begin{obligation}[Uniform construction from a three-counter machine]
\label{obl:construction}
Given a finite symmetric three-counter machine \(M\), using only
effective finite constructions, Presburger operations, and the
algorithms S2 and S3, construct the following finite data and prove
the listed properties.
\begin{enumerate}[label=\textnormal{(C\arabic*)},leftmargin=*]
\item \textbf{Representations.}
Construct a finite union \(X\) of labelled counter spaces of
dimensions at most three, a finite scalar space \(A=Q_A\times\N\),
and an explicit Presburger encoding
\[
 \jmath:\operatorname{Conf}(M)\longrightarrow X.
\]
For example, each original control can have a designated
three-dimensional copy with \(\jmath(q,x)=(\widehat q,x)\).
Specify the physical meanings of all nonprivate representations.

\item \textbf{The relations defining \(K\).}
Specify \(Z_0,R_{\mathrm{in}},R_{\mathrm{out}},B_0,E\) with the
types of Definition~\ref{def:K-data}, by finite graphs of actual
run relations and explicit representation changes.
The graphs, permitted edge types, and start/end labels must all be
fixed independently of any endpoint box query.

\item \textbf{The scalar and query properties.}
Supply the algorithms in (K1)--(K4) of
Theorem~\ref{thm:K-algorithm}. In particular, prove pumping (LP)
for the particular relation \(E\).
For (K4), supply an exact two-counter presentation, such as that in
Lemma~\ref{lem:band-compiler}, or another proved construction.
For every use of Theorem~\ref{thm:face-query}, exhibit its finite
graph and verify (F1)--(F4), including ordinary-cycle weights and
weak monotonicity.

\item \textbf{Independent events.}
Supply a common period \(D>0\) and finite event lists defining
\(U,V\) exactly as in Definition~\ref{def:independent-events},
including the effective eligibility bases.
Private events must have the graph restrictions and absorption
identities needed by Lemma~\ref{lem:private-pump}.

\item \textbf{Exact correspondence with the original problem.}
For the fixed relation \(K\) of \eqref{eq:defined-K}, prove
\begin{equation}
 c\longrightarrow_M^*d
 \quad\Longleftrightarrow\quad
 \jmath(c)\,
 K;\bigl((U\cup V);K\bigr)^*
 \,\jmath(d)
 \qquad(c,d\in\operatorname{Conf}(M)).
 \label{eq:machine-equivalence}
\end{equation}
The forward implication must account for every original run.
The reverse implication must account for every allowed composition
of the supplied relations, including private blocks.
\end{enumerate}
\end{obligation}

The statement asks for one presentation uniform in all endpoint pairs
of \(M\). For point-to-point decidability, a weaker version also
suffices: construct the data effectively from \((M,c,d)\), require
\eqref{eq:machine-equivalence} only for that pair, and keep the resulting
relations fixed throughout every internal box query. The distinction
does not affect the algorithm once its data are supplied.

This is the point at which all local arguments must agree.
Constructing a family for which the pair bound holds, and separately
constructing a larger family into which all runs normalize, is
insufficient unless they are proved to be the same family.
Likewise, an algorithm for literal event-avoiding paths does not
automatically decide a larger \(K\) containing extra replacement
words.

\begin{theorem}[The proved implication from the construction]
\label{thm:conditional-conclusion}
If Construction obligation~\ref{obl:construction} is fulfilled
uniformly, then under S2 and S3, reachability of symmetric
three-counter machines is decidable.
\end{theorem}
\begin{proof}
Construct the data for the given machine.
Theorem~\ref{thm:K-algorithm} gives a box-intersection algorithm for
the fixed \(K\).
Theorem~\ref{thm:formal-completion} decides
\[
 \Hit_{\,K;((U\cup V);K)^*}
       \bigl(\{\jmath(c)\},\{\jmath(d)\}\bigr),
\]
since singletons are grid boxes.
By \eqref{eq:machine-equivalence}, this answer is exactly whether
\(c\longrightarrow_M^*d\).
\end{proof}

\subsection{What the earlier text did not justify}

The earlier version stated that the normalization supplied one fixed
relation \(K\) with a box algorithm, but its descriptions did not
fully specify and verify the generated family required in
(C2)--(C5). The decreasing measure proved termination only
\emph{assuming} appropriate replacements. It did not prove that each
replacement exists effectively, stays in the same specified update
graphs, and preserves all conditions used by the box algorithm.
Consequently the earlier assertion that the S2--S3 decidability proof
was complete was too strong.

The following obligations are particularly concrete.
\begin{itemize}[leftmargin=*]
\item For each Cartesian or axis replacement, give the finite
      list of relations in (PK) or (VL), prove their inclusions,
      and prove that they remain in the family used in (C2).
\item For each transfer record, prove (TR) at its named anchor.
      If an expansion is followed by a negative relation, a bound
      on the canonical representative's pair cannot replace a
      bound on the actual expanded vector.
      The proposed private negative event must be justified by
      its particular identity \(V;R=R\).
\item For the ordinary graph, verify the absence of negative
      ordinary cycles \emph{after} all required finite word
      implementations have been included.
      The constants \(m,P,R,B\) used to select the large-gain cutoff
      must already be bounded independently of that cutoff.
\item For each first/last scalar relation, give a formula-producing
      algorithm as required in (K2).
      A membership test, an emptiness test, or a claim of
      semilinearity without an effective construction does not
      fulfill this requirement.
\item For the finite scalar relation \(E\), prove (LP) with a
      common computable period while retaining its endpoint
      encodings and all path restrictions.
      S3 and S2 supply formulas for finite compositions; they do
      not imply this pumping statement for arbitrary such formulas.
\end{itemize}

No counterexample to the desired construction is asserted here.
The distinction is between a fully verified construction and the
local arguments that would support one. The target theorem is not
declared false or known to be open in the literature; this document
records the present proof status.

\subsection{A precise checkpoint for continuation}

The following parts are proved in this version:
\begin{enumerate}[leftmargin=*]
\item the exact S3 boundary reduction,
      Proposition~\ref{prop:boundary};
\item effective upward-basis extraction and the event-completion
      algorithm, Theorem~\ref{thm:formal-completion};
\item scalar saturation under its stated pumping hypotheses,
      Theorem~\ref{thm:formal-scalar};
\item the pair-growth and counting lemmas,
      Lemmas~\ref{lem:cycle-pair-growth} and
      \ref{lem:relation-count};
\item the box-query algorithm for a supplied face graph,
      Theorem~\ref{thm:face-query};
\item the bounded-band two-counter compiler,
      Lemma~\ref{lem:band-compiler};
\item the exact fixed-\(K\) formula and its algorithmic consequence,
      Theorem~\ref{thm:K-algorithm}.
\end{enumerate}
The next research task is to fulfill
Construction obligation~\ref{obl:construction}, starting with the
actual finite vertex and edge lists for (C2), then proving the
replacement properties for those lists. A later proof must refer to
those fixed relations throughout; a new name for a path family does
not remove the need to verify its properties.

\paragraph{Supporting record.}
The earlier geometric attempts, local arguments, and failed approaches
remain in the
\href{https://chatgpt.com/library?file_id=libfile_f10d80492b688191970c625ab3952096&account_id=personal}{research log}.
The
\href{https://chatgpt.com/library?file_id=libfile_7338a158ba3c8191a7b47b1688ce2dd5&account_id=personal}{earlier assembled text}
is retained as a historical version. Its statements of completion
must be read with the correction in this section.

\end{document}
